% =====================================================================
%  LA CRYPTE DE L'ORDRE DES NOMBRES PREMIERS
%  Sujet ludique — Mathématiques expertes, terminale générale
%  Quête de calcul matriciel : somme, différence, produit, inversion,
%  comparaison de matrices 2x2, 3x3, 4x4 — chiffrement modulo 7 et 13.
%  Compilation : pdflatex (2 passes) ou Overleaf.
% =====================================================================
\documentclass[11pt,a4paper]{article}

\usepackage[utf8]{inputenc}
\usepackage[T1]{fontenc}
\usepackage[french]{babel}
\usepackage{amsmath,amssymb}
\usepackage[margin=2.1cm,bottom=2.6cm]{geometry}
\usepackage{booktabs,array}
\usepackage{enumitem}
\usepackage{xcolor}
\usepackage{tcolorbox}
\usepackage{fancyhdr}


% ---------- Couleurs de la quête ----------
\definecolor{encre}{RGB}{26,32,44}
\definecolor{or}{RGB}{153,110,20}
\definecolor{argent}{RGB}{95,112,130}
\definecolor{emeraude}{RGB}{4,120,87}
\definecolor{jade}{RGB}{10,90,74}
\definecolor{parche}{RGB}{253,249,236}
\definecolor{nuit}{RGB}{24,28,44}
\definecolor{sang}{RGB}{150,32,32}

% ---------- Boîtes ----------
\newtcolorbox{missive}[1]{colback=parche,colframe=or,title={\textbf{#1}},
  coltitle=white,fonttitle=\normalsize,boxrule=1pt,arc=2mm,left=3mm,right=3mm}
\newtcolorbox{rituel}[1]{colback=white,colframe=argent,title={\textbf{#1}},
  coltitle=white,fonttitle=\normalsize,boxrule=1pt,arc=2mm,left=3mm,right=3mm}
\newtcolorbox{piege}{colback=sang!4,colframe=sang!80,boxrule=0.8pt,arc=1mm,
  left=3mm,right=3mm}
\newtcolorbox{conseil}{colback=emeraude!5,colframe=emeraude!70,boxrule=0.8pt,
  arc=1mm,left=3mm,right=3mm}
\newtcolorbox{precedemment}[1]{colback=argent!6,colframe=nuit!75,
  title={\textbf{#1}},coltitle=white,fonttitle=\normalsize,boxrule=1pt,
  arc=2mm,left=3mm,right=3mm}

% ---------- En-têtes ----------
\pagestyle{fancy}
\fancyhf{}
\fancyhead[L]{\small\textsc{La Crypte de l'Ordre des Nombres Premiers}}
\fancyhead[R]{\small\textit{Maths expertes, terminale générale}}
\fancyfoot[L]{\small\textit{Quête de calcul matriciel}}
\fancyfoot[C]{\small\thepage}
\fancyfoot[R]{\small\textit{Récompense de la quête : /20}}
\renewcommand{\headrulewidth}{0.6pt}

\setlength{\parindent}{0pt}
\setlist{itemsep=2pt,topsep=3pt}

% ---------- Raccourcis ----------
\newcommand{\Z}{\mathbb{Z}}
\newcommand{\Id}[1]{I_{#1}}
\newcommand{\Epx}[1]{\textcolor{or}{\textbf{[#1 XP]}}}

\begin{document}

% =====================================================================
\begin{center}
  {\LARGE\bfseries\textcolor{nuit}{LA CRYPTE DE L'ORDRE DES NOMBRES PREMIERS}}\\[2mm]
  {\large\itshape Une quête de calcul matriciel pour apprenti·e cryptographe}\\[2mm]
  \textbf{Mathématiques expertes, terminale générale}\quad\textbullet\quad
  Feuilleton en deux épisodes de 50 min\quad\textbullet\quad Barème : 20 XP (10 + 10)
\end{center}

\vspace{1mm}
\begin{missive}{Message retrouvé dans le tombeau du dernier Grand Maître}
\itshape
Apprenti·e, si tu lis ces lignes, la Crypte t'a accepté·e.
Des siècles durant, notre Ordre a protégé son trésor derrière trois sceaux
de matrices. Les nombres premiers en sont l'alphabet, et les restes,
seuls, disent la vérité : \textbf{seul le reste parle}.

Trois épreuves t'attendent, de plus en plus profondes :
le \textcolor{argent}{\textbf{Sceau d'argent}} (matrices $2\times 2$, rituel modulo 7),
le \textcolor{emeraude}{\textbf{Sceau d'émeraude}} (matrices $3\times 3$, rituel modulo 13)
et le \textcolor{jade}{\textbf{Sanctuaire de jade}} (matrices $4\times 4$, rituel modulo 13).

Chaque épreuve livrera un fragment du message perdu de l'Ordre.
Assemble-les, et le nom du trésor t'apparaîtra. Mais n'espère pas tout
traverser en un jour : la Crypte ne s'ouvre que deux fois, cinquante
minutes à chaque lune, et elle se referme entre les deux plongées.
Le coffre final ne s'ouvrira qu'avec le code des premiers.
Que la quête commence !
\end{missive}

% =====================================================================
\section*{\textcolor{nuit}{Générique : règles de la quête}}

{\small\itshape (Le générique passe avant chaque épisode : ces rituels
valent pour toute la saison.)}

\begin{rituel}{Les quatre rituels de l'Ordre}
\begin{enumerate}[label=\textbf{R\arabic*.},leftmargin=10mm]
  \item \textbf{Seul le reste parle.} Dans la Crypte, tout est gravé en
        \emph{nombres premiers}. Chaque nombre premier $p$ gravé sur la pierre
        doit d'abord être réduit : on ne retient que son reste dans la division
        euclidienne par le module du sceau (7 ou 13).
        Par exemple, dans le sanctuaire (module 13) : $61 \equiv 9 \pmod{13}$,
        car $61 = 4\times 13 + 9$.
  \item \textbf{Les matrices vivent modulo $n$.} Deux matrices sont
        \emph{égales dans le sceau} si leurs coefficients sont égaux modulo $n$
        (7 ou 13). Tous les produits de matrices se réduisent modulo $n$.
  \item \textbf{L'alphabet du chiffre.} Les messages sont écrits en lettres
        majuscules non accentuées (l'Ordre a perdu ses accents : É se lit E).
        \begin{itemize}
          \item Dans les sceaux \textbf{modulo 13}, chaque nombre $x$
                ($0\leqslant x \leqslant 12$) code une lettre :
                $\text{A}=0,\ \text{B}=1,\ \dots,\ \text{M}=12$
                (alphabet réduit de l'Ordre, annexe B).
          \item Dans le sceau \textbf{modulo 7}, chaque lettre est portée par un
                \emph{couple} $(a,b)$ de nombres ($0\leqslant a,b\leqslant 6$) :
                le rang de la lettre ($\text{A}=0$, $\dots$, $\text{Z}=25$) vaut
                $a + 7b$ (annexe C).
        \end{itemize}
  \item \textbf{La Table de l'Ordre.} Pour ouvrir le coffre final, chaque lettre
        vaut un nombre premier selon la Table de l'Ordre (annexe A) :
        $\text{A}=2,\ \text{B}=3,\ \text{C}=5$, etc.
\end{enumerate}
\end{rituel}

\begin{conseil}
\textbf{Mémo du cryptographe} (annexe D) : inverse d'une matrice $2\times 2$,
produit de matrices, et le grand théorème des voleurs de clés :
$(M N)^{-1} = N^{-1} M^{-1}$. Les annexes restent à ta disposition
lors des deux plongées.
\end{conseil}

% =====================================================================
\section*{\textcolor{nuit}{Épisode 1 : La Pierre de Voûte}}
{\small\textbf{Première séance : 50 min}\quad\textbullet\quad
\textbf{Récompense : 10 XP}\quad\textbullet\quad
Acte I : rituel d'entrée 10 min\quad\textbullet\quad Acte II : le Sceau
d'argent 25 min\quad\textbullet\quad Acte III : le parchemin du traître
15 min}

{\small\itshape (Acte I : lis le générique et jette un œil aux annexes.
Ensuite, la crypte t'appartient... jusqu'au bout des cinquante minutes.)}

\subsection*{\textcolor{argent}{Acte II : Le Sceau d'argent (épreuve 1)} \normalsize\textcolor{argent}{(niveau 1 : matrices $2\times 2$, rituel modulo 7, 25 min)} : 6 XP}

La Porte d'argent garde la première pierre. Sur elle sont gravées trois
matrices :
\[
A=\begin{pmatrix}3&1\\2&5\end{pmatrix}\qquad
B=\begin{pmatrix}4&6\\0&4\end{pmatrix}\qquad
C=\begin{pmatrix}5&4\\1&5\end{pmatrix}.
\]

\begin{enumerate}[label=\textbf{A.\arabic*.},leftmargin=10mm]
  \item \textbf{La pierre de voûte.} Calculer la matrice
        $K_1 = A + B - C$.\Epx{1}
  \item \textbf{Le faux gravage.} Un vandale a prétendu que la pierre de voûte
        était l'une de ces trois matrices :
        \[
        S_1=\begin{pmatrix}2&3\\1&4\end{pmatrix}\qquad
        S_2=\begin{pmatrix}3&2\\4&1\end{pmatrix}\qquad
        S_3=\begin{pmatrix}2&-3\\1&4\end{pmatrix}.
        \]
        \textbf{Comparer} $K_1$ avec $S_1$, $S_2$ et $S_3$ : quelle est la vraie
        pierre ?\Epx{0,5}
  \item \textbf{La clé s'inverse.} On travaille désormais modulo 7.
        \begin{enumerate}[label=\alph*)]
          \item Calculer $\det K_1$, puis vérifier que $3\times 5 - 2\times 7 = 1$.
                En déduire l'inverse de $5$ modulo 7.\Epx{0,5}
          \item Montrer que
                $K_1^{-1} \equiv \begin{pmatrix}5&5\\4&6\end{pmatrix} \pmod 7$,
                puis \textbf{vérifier} que $K_1 K_1^{-1} \equiv I_2 \pmod 7$.\Epx{1,5}
        \end{enumerate}
  \item \textbf{Le premier fragment.} Plus profond dans la crypte, un
        linteau porte quatre nombres premiers gravés :
        \[
        \begin{pmatrix}11 & 7\\ 29 & 7\end{pmatrix}.
        \]
        \begin{enumerate}[label=\alph*)]
          \item Appliquer le rituel R1 (module 7) pour obtenir la matrice
                chiffrée $Y$.\Epx{0,5}
          \item Calculer $X = K_1^{-1}Y \pmod 7$.\Epx{1}
          \item Chaque colonne $(a,b)$ de $X$ porte une lettre de rang $a+7b$
                (rituel R3). \textbf{Décrypter} le premier fragment du message.\Epx{1}
        \end{enumerate}
\end{enumerate}

% =====================================================================
\subsection*{\textcolor{emeraude}{Acte III : Le Sceau d'émeraude (épreuve 2, premier affrontement)} \normalsize\textcolor{emeraude}{(niveau 2 : matrices $3\times 3$, rituel modulo 13, 15 min)} : 4 XP}

L'écho de la Porte d'argent s'éteint derrière toi. Au fond du couloir,
une lueur verte pulse : le Sceau d'émeraude. Il fut scellé par deux
frères : l'un gravait les lignes du haut, l'autre celles du bas :
\[
M=\begin{pmatrix}1&2&0\\0&1&3\\0&0&1\end{pmatrix}\qquad
N=\begin{pmatrix}1&0&0\\4&1&0\\0&5&1\end{pmatrix}.
\]

\begin{enumerate}[label=\textbf{B.\arabic*.},leftmargin=10mm]
  \item \textbf{Le sceau composé.} Calculer le produit $M\times N$, puis le
        réduire modulo 13 : on obtient le sceau $K_2$.\Epx{2}
  \item \textbf{Le parchemin du traître.} Un traître a laissé deux parchemins,
        prétendant chacun contenir la clé de décryptage $K_2^{-1}$.
        Les archives de l'Ordre fournissent les matrices
        \[
        M'=\begin{pmatrix}1&11&6\\0&1&10\\0&0&1\end{pmatrix}\quad\text{et}\quad
        N'=\begin{pmatrix}1&0&0\\9&1&0\\7&8&1\end{pmatrix},
        \]
        et l'on sait que $M M' = I_3$ et $N N' = I_3$ modulo 13
        (à vérifier si tu doutes de l'Ordre !).
        \begin{enumerate}[label=\alph*)]
          \item Calculer les deux produits $P_1 = M'N'$ et $P_2 = N'M'$
                modulo 13. Sont-ils égaux ?\Epx{1}
          \item \textbf{Comparer} : calculer $K_2 P_1$ et $K_2 P_2$ modulo 13.
                Lequel des deux parchemins est le faux ? Quel est le vrai
                $K_2^{-1}$ ?\Epx{1}
        \end{enumerate}
\end{enumerate}

\begin{piege}
\textbf{Avertissement du Gardien :} le traître croyait que l'ordre des
facteurs ne comptait pas. Mais l'inverse d'un produit renverse l'ordre :
$(MN)^{-1} = N^{-1}M^{-1}$, et non $M^{-1}N^{-1}$. Les deux parchemins ne
pouvaient être vrais à la fois... et pourtant, l'un des deux l'est bel et
bien !
\end{piege}

% ---------------------------------------------------------------------
\begin{missive}{La crypte se referme... (fin de l'épisode 1)}
\itshape
Le Gardien lève sa torche : le temps de la première plongée est écoulé.
Derrière toi, la Porte d'argent s'est refermée à jamais ; devant toi, le
Sceau d'émeraude reste \textbf{muet}. Tu ne possèdes qu'un fragment du
message : \textbf{« LA »}. La quoi ? La lumière défaille... En sortant,
tu remarques une inscription gravée au sol :
\upshape
\begin{center}
« \textbf{La clé que tu cherches porte un nom de pierre.} »
\end{center}
\itshape
Avant que la crypte ne se referme, vérifie ton carnet : il doit contenir
$K_1$, $K_1^{-1}$, $K_2$ et surtout le vrai parchemin $K_2^{-1}$.
Le Gardien ne les redonnera pas : ils te serviront à la prochaine lune.
\upshape

\textbf{Prochain épisode :} un mot de trois lettres gravé dans
l'émeraude, trois gardiens menteurs au fond du sanctuaire, et un coffre
qui ne s'ouvre qu'avec le code des nombres premiers...
\end{missive}

\begin{center}
\itshape (Fin de l'épisode 1. L'épisode 2, à la prochaine lune...)
\end{center}

	% =====================================================================
\newpage
\section*{Annexes de la Crypte}

\subsection*{Annexe A : La Table de l'Ordre (nombres premiers)}
\begin{center}
	\begin{tabular}{*{13}{c}}
		\toprule
		A & B & C & D & E & F & G & H & I & J & K & L & M\\
		\midrule
		2 & 3 & 5 & 7 & 11 & 13 & 17 & 19 & 23 & 29 & 31 & 37 & 41\\
		\midrule
		N & O & P & Q & R & S & T & U & V & W & X & Y & Z\\
		\midrule
		43 & 47 & 53 & 59 & 61 & 67 & 71 & 73 & 79 & 83 & 89 & 97 & 101\\
		\bottomrule
	\end{tabular}
\end{center}

\subsection*{Annexe B : Alphabet réduit (sceaux modulo 13)}
Chaque nombre $x$ de $0$ à $12$ code une lettre :
\begin{center}
	\begin{tabular}{l|*{13}{c}}
		& x=0 & 1 & 2 & 3 & 4 & 5 & 6 & 7 & 8 & 9 & 10 & 11 & 12\\
		\hline
		Lettre & A & B & C & D & E & F & G & H & I & J & K & L & M\\
	\end{tabular}
\end{center}

\subsection*{Annexe C : Alphabet complet par paires (sceau modulo 7)}
Chaque lettre de rang $\ell$ ($\text{A}=0$, $\dots$, $\text{Z}=25$) est portée
par le couple $(a,b)$ tel que $\ell = a + 7b$ :
\begin{center}
	\begin{tabular}{*{13}{c}}
		\toprule
		A & B & C & D & E & F & G & H & I & J & K & L & M\\
		(0,0) & (1,0) & (2,0) & (3,0) & (4,0) & (5,0) & (6,0) & (0,1) & (1,1) & (2,1) & (3,1) & (4,1) & (5,1)\\
		\midrule
		N & O & P & Q & R & S & T & U & V & W & X & Y & Z\\
		(6,1) & (0,2) & (1,2) & (2,2) & (3,2) & (4,2) & (5,2) & (6,2) & (0,3) & (1,3) & (2,3) & (3,3) & (4,3)\\
		\bottomrule
	\end{tabular}
\end{center}

\subsection*{Annexe D : Mémo du cryptographe}
\begin{itemize}
	\item \textbf{Produit :} si $M$ ($n$ lignes, $p$ colonnes) et $N$
	($p$ lignes, $q$ colonnes), le coefficient $(i,j)$ de $MN$ est
	$\sum_{k=1}^{p} m_{ik}\, n_{kj}$.
	\item \textbf{Inverse $2\times 2$ :} si $D = \begin{pmatrix}a&b\\c&d\end{pmatrix}$
	avec $\delta = ad-bc \neq 0$, alors
	$D^{-1} = \dfrac{1}{\delta}\begin{pmatrix}d&-b\\-c&a\end{pmatrix}$.
	Modulo $n$ : remplacer $\frac{1}{\delta}$ par l'inverse de $\delta$
	modulo $n$ (s'il existe).
	\item \textbf{Le théorème des voleurs de clés :}
	$(MN)^{-1} = N^{-1}M^{-1}$ : l'ordre s'inverse.
	\item \textbf{Inversibles modulo $n$ :} $x$ est inversible modulo $n$
	si et seulement si $x$ et $n$ sont premiers entre eux
	(Bézout). C'est pourquoi l'Ordre chérit les modules
	premiers : tout $x$ non nul y est inversible.
\end{itemize}

\end{document}
